\chapter{Securite}

\section{Gestion des secrets --- SSM Parameter Store}

Les secrets sont stockes dans AWS Systems Manager Parameter Store
en tant que \texttt{SecureString} (chiffres avec KMS)~:

\begin{table}[H]
    \centering
    \begin{tabularx}{\textwidth}{lXl}
        \toprule
        \textbf{Parametre} & \textbf{Usage} & \textbf{Type} \\
        \midrule
        /deus-dashboard/dev/nasa-api-key & Cle API NASA & SecureString \\
        /deus-dashboard/dev/pg-host & Adresse RDS & String \\
        /deus-dashboard/dev/pg-password & Mot de passe DB & SecureString \\
        /deus-dashboard/dev/dynamodb-table & Nom de la table & String \\
        \bottomrule
    \end{tabularx}
    \caption{Parametres SSM}
\end{table}

Le role d'execution ECS lit ces parametres au demarrage du conteneur
et les injecte comme variables d'environnement.

\section{IAM --- Adaptations Learner Lab}

\subsection{Architecture cible vs. realite Learner Lab}

L'architecture ideale prevoyait trois roles IAM distincts suivant le
principe du moindre privilege~:

\begin{description}
    \item[Execution Role] Permettrait a ECS de~: tirer des images ECR,
        lire les parametres SSM, ecrire dans CloudWatch Logs.
    \item[Task Role] Permettrait au conteneur de~: acceder a DynamoDB
        (lecture/ecriture sur la table de cache).
    \item[Instance Role] Permettrait aux instances EC2 de~: s'enregistrer
        au cluster ECS, utiliser SSM Session Manager, envoyer des metriques CloudWatch.
\end{description}

\subsection{Contrainte Learner Lab}

L'environnement AWS Academy Learner Lab interdit la creation de roles IAM
(\texttt{iam:CreateRole} non autorise). En consequence, l'ensemble des
ressources utilisent le role pre-existant \textbf{LabRole}, qui dispose
de permissions larges couvrant les services necessaires (ECS, ECR, SSM,
CloudWatch, DynamoDB, S3, RDS).

\begin{lstlisting}[language=bash,caption={Utilisation du LabRole dans Terraform}]
# Au lieu de creer des roles personnalises
data "aws_iam_role" "lab_role" {
  name = "LabRole"
}

# Execution role et task role ECS
execution_role_arn = data.aws_iam_role.lab_role.arn
task_role_arn      = data.aws_iam_role.lab_role.arn

# Instance profile pour les EC2 (ECS, bastion)
data "aws_iam_instance_profile" "lab_profile" {
  name = "LabInstanceProfile"
}
\end{lstlisting}

\textbf{Compromis de securite~:} Le \texttt{LabRole} dispose de permissions
plus larges que necessaire, ce qui contrevient au principe du moindre privilege.
En environnement de production, on creerait des roles specifiques avec des
politiques restrictives pour chaque composant.

\section{Reseau --- Segmentation VPC}

\begin{table}[H]
    \centering
    \begin{tabularx}{\textwidth}{lXl}
        \toprule
        \textbf{Security Group} & \textbf{Regles entrantes} & \textbf{Cible} \\
        \midrule
        sg-alb & 80 depuis 0.0.0.0/0 & ALB \\
        sg-ecs & 32768-65535 depuis sg-alb, 22 depuis sg-bastion & Instances ECS \\
        sg-rds & 5432 depuis sg-ecs et sg-bastion & PostgreSQL \\
        sg-bastion & 22 depuis IP autorisee & Bastion \\
        sg-monitoring & 3000/9090 depuis IP autorisee, 22 depuis sg-bastion & Prometheus/Grafana \\
        \bottomrule
    \end{tabularx}
    \caption{Security Groups et regles}
\end{table}

\textbf{Moindre privilege par references de Security Groups~:} chaque flux
interne est autorise par \emph{reference de Security Group} et non par plage
CIDR. Initialement, la base RDS acceptait le port 5432 depuis l'ensemble des
sous-reseaux prives~; la regle a ete restreinte aux seuls
\texttt{sg-ecs} et \texttt{sg-bastion}. Ainsi, meme une ressource compromise
placee dans un sous-reseau prive ne peut pas atteindre la base si elle ne
porte pas le bon Security Group. De meme, le port 443 de l'ALB, sans
listener HTTPS associe (certificat ACM indisponible en Learner Lab),
a ete supprime~: aucune regle ne reste ouverte sans service derriere.

Les flux autorises se resument a~:
\begin{lstlisting}[language={},caption={Flux reseau (moindre privilege)}]
Internet --80--> sg-alb --ports dynamiques--> sg-ecs --5432--> sg-rds
Admin    --22--> sg-bastion --22--> sg-ecs / sg-monitoring
                 sg-bastion --5432--> sg-rds   (administration DB)
Admin    --3000/9090--> sg-monitoring          (UI Grafana/Prometheus)
\end{lstlisting}

Seul l'ALB est expose sur Internet. Les instances ECS et la base de donnees
sont dans des sous-reseaux prives, accessibles uniquement via l'ALB
ou le bastion SSH.

\section{Bastion SSH}

Le bastion est le point d'entree unique pour l'acces SSH aux ressources privees.
Il est configure avec~:
\begin{itemize}
    \item Acces SSH restreint a l'IP de l'administrateur
    \item Instance \texttt{t2.micro} (Free Tier)
    \item Durcissement via Ansible (pas de root SSH, pas d'authentification par mot de passe)
\end{itemize}

\section{Cle API NASA}

La cle API NASA, auparavant codee en dur dans le code source Angular
(visible dans le navigateur), est desormais~:
\begin{enumerate}
    \item Stockee dans SSM Parameter Store (chiffree)
    \item Injectee dans le conteneur backend via la Task Definition ECS
    \item Utilisee cote serveur uniquement --- le frontend ne la voit plus
\end{enumerate}
